\section{Related work}

The closest comparisons concern the population contrast between assigning treatment to everyone and assigning treatment to no one. For this target, precision depends jointly on the response class, the assignment mechanism, and the network information supplied to the estimator. The present analysis combines additive responses with bounded coefficient mass and bounded in- and out-degrees, independent treatment assignment with probability one half, and independent edge retention with a known probability. Over the schedule class in \cref{obj:def:model}, the risk in \cref{obj:def:risk} allows every measurable randomized estimator to use the retained graph, the full assignment vector, and all realized outcomes. The characterization in \cref{obj:thm:observed-record-precision-frontier} therefore describes optimal squared-error precision for this complete observation experiment.

Known-neighborhood SNIPE provides the closest estimator comparison. It targets the same population contrast under independent Bernoulli assignment and polynomial neighborhood responses of bounded interaction order. Its published analysis establishes unbiasedness, a variance upper bound, and a minimax lower bound \citep[Sections 3--5; Theorems 1--2]{CortezRodriguezEichhornYu2023SNIPE}. The present additive model corresponds to interaction order one. Because the published neighborhoods include each recipient itself, the common off-diagonal degree bound \leanref{sym:d}{\ensuremath{d}} in \cref{obj:def:model} translates into augmented in- and out-degree bounds of \(d+1\); the coefficient-mass normalization is preserved by this embedding, as recorded in \cref{obj:lem:published-model}. For population size \leanref{sym:n}{\ensuremath{n}}, the published variance bound specializes to order \((d+1)^3/n\) at treatment probability one half. Its minimax lower bound uses Gaussian coefficient priors and, at interaction order one, a common treatment probability below \(0.16\). These statements furnish distinct comparisons for the bounded additive class and assignment mechanism studied here. Under full edge retention, \cref{obj:thm:observable-upper} supplies a quadratic degree upper bound. The sharper supplied-graph benchmark is stated in \cref{obj:thm:supported-boundaries}.
% [SNIPE source](https://arxiv.org/pdf/2208.05553)

A second comparison concerns optimization of the experiment itself. With the interference network supplied, design-based minimax analysis can optimize jointly over assignment mechanisms and estimators under an arbitrary neighborhood response model and its specified potential-outcome envelope. For random regular networks whose degree grows more slowly than the square root of population size, the global average treatment effect has minimax risk bounded below by degree divided by population size and above by squared degree divided by population size, with high probability over the network \citep[Corollary 5.4]{KandirosHarshawSavje2026DesignBasedMinimax}. That comparison concerns a fixed supplied network and design optimization under its own response envelope. Here, the assignment law is fixed, the response envelope bounds absolute additive coefficient mass, and the worst case ranges over graphs and coefficients. The supplied-graph risk in \cref{obj:def:supplied-risk} isolates the information comparison while retaining the paper's assignment mechanism and schedule class.
% [Design-based minimax source](https://arxiv.org/html/2608.04909v1)

Network measurement error also motivates correction of exposure classifications. Under a four-level exposure model and Bernoulli treatment assignment, known false-positive and false-negative edge probabilities support method-of-moments correction of exposure-specific mean outcomes. The resulting guarantees include asymptotic bias and variance control and, under additional conditions, asymptotic normality \citep[Section 3; Theorems 4--7]{LiSussmanKolaczyk2021NetworkNoise}. This work credits the role of known network-error probabilities in correcting causal estimators. The present observation model instead supplies a randomly retained subgraph of the true directed graph, and the response class permits heterogeneous additive effects from individual sources. Its precision guarantee concerns the all-treated-versus-all-control contrast and worst-case squared loss over all estimators.
% [Network-noise source](https://arxiv.org/pdf/2105.04518)

Supplied supergraphs offer another form of partial network information. UNITE estimates the global average treatment effect under Bernoulli assignment using a supplied neighborhood superset. Its additive and bounded-order polynomial analyses establish unbiasedness and consistency under the stated containment and response conditions, together with variance bounds whose constants depend on neighborhood sizes, assignment probabilities, and the response envelope \citep[Theorems 1, 4, and 8; Appendix B.1]{ShankarEtAl2024UNITE}. A containing supergraph supplies every relevant source as a candidate, whereas a retained subgraph supplies sampled true arrows together with their inclusion probability. Baseline information provides a further route: under heterogeneous additive responses and equal marginal treatment probabilities along interfering edges, an observed baseline supports an unbiased graph-free estimator of the total treatment effect \citep[Corollary 3]{YuAiroldiBorgsChayes2022}. These results clarify how neighborhood containment and baseline observations change the estimator's information set. In the present experiment, baselines are fixed schedule coefficients, and estimation uses the retained graph, assignments, and realized outcomes.
% [UNITE source](https://proceedings.mlr.press/v238/shankar24a/shankar24a.pdf), [Observed-baseline source](https://pmc.ncbi.nlm.nih.gov/articles/PMC9636977/)

Approximate-network analysis studies a related bias--variance trade-off for a specified surrogate-network estimator. Under bounded outcomes and restrictions on first- and second-order treatment influences, its asymptotic variance is \(O\{d_{\mathcal G}^2/[n\pi(1-\pi)]\}\), where \(d_{\mathcal G}\) is surrogate-network degree and \(\pi\) is treatment probability; constant outcomes on a regular surrogate graph give the same estimator-specific order. A separate omitted-influence condition controls relative bias \citep[Assumptions 1--3; Theorems 1--2; Lemma 2]{JiangDengWangWang2024ApproximateNetworks}. Those guarantees describe the performance of a chosen estimator as the surrogate network captures different amounts of influence. The present characterization optimizes squared risk over the full estimator class under the explicit edge-retention experiment. Its lower bound applies to unrestricted use of the observed record, while its upper bound supplies an attainable procedure under the same schedule restrictions. This common experiment makes the contribution interpretable as a precision characterization indexed by population size, degree, and edge retention.
% [Approximate-network source](https://arxiv.org/html/2408.04441v1)
